% GENERATED FILE. Edit canonical lessons, then run npm run generate:books.
% canonical-source-hash: fnv1a64:fdc2c223a97f7cf1
% canonical-lessons: TA-C497-koncu, TA-C497-cirattu, TA-C497-kilppati, TA-C497-paramari, TA-C497-alocanai

\chapter{Raising Children}
\label{ch:ta-raising-children}
\hlchaptermodality{497}

\begin{chapteropening}
\textbf{By the end of this chapter:} \emph{I can talk about doting on children, pampering, obedience, care and advice.}

Complete the last lesson of chapter 497: I can talk about doting on children, pampering, obedience, care and advice.
\end{chapteropening}

\section[\texorpdfstring{\ta{கொஞ்சு} (koñcu)}{கொஞ்சு (koñcu)}]{\ta{கொஞ்சு} (koñcu) — to coo over; to fondle}

\label{lesson:TA-C497-koncu}



\begin{quote}
Before the new one: say the Tamil for a son or daughter, then the Tamil for to adopt.
\end{quote}

\subsection*{You'll want to know: \ta{கொஞ்சு}}
\textbf{\ta{கொஞ்சு}} — \emph{koñcu} — \textquotedblleft{}to coo over; to fondle\textquotedblright{}.

\begin{quote}
\ta{தாத்தா} \ta{பேத்தியைக்} \ta{கொஞ்சினார்}.
\end{quote}

\begin{quote}
\emph{tāttā pēttiyaik koñciṉār.}
\end{quote}

\begin{quote}
Grandfather cooed over his granddaughter.
\end{quote}

Past (I …): \ta{கொஞ்சினேன்}, \emph{koñciṉēṉ}.

Also of a child’s sweet prattle.

\begin{cousinweb}
\textbf{\ta{கொஞ்சு}} (\emph{koñcu}) is native Dravidian (DEDR 2043), for loving baby talk and fondling.
\end{cousinweb}

\begin{grammarlens}[title={What you've built}]
One more word for describing what you lived and felt.
\end{grammarlens}

\subsection*{Your turn}
\begin{itemize}
\raggedright
  \item \emph{Say it:} \emph{koñcu}
  \item \emph{Say it:} \emph{koñcu}, once more
  \item \emph{Say it:} \emph{tāttā pēttiyaik koñciṉār.} — Grandfather cooed over his granddaughter.
  \item \emph{Recall:} say the Tamil for a son or daughter, then the Tamil for to adopt, then say \emph{koñcu} again
\end{itemize}

\begin{tcolorbox}[breakable,colback=teal!4,colframe=teal!35!black,title={Before you move on}]
What does \ta{கொஞ்சு} mean? (\textquotedblleft{}To coo over; to fondle\textquotedblright{}.) Say it once more.
\end{tcolorbox}
\section[\texorpdfstring{\ta{சீராட்டு} (cīrāṭṭu)}{சீராட்டு (cīrāṭṭu)}]{\ta{சீராட்டு} (cīrāṭṭu) — to pamper; to cherish}

\label{lesson:TA-C497-cirattu}



\begin{quote}
Before the new one: say the Tamil for to adopt, then the Tamil for to coo over.
\end{quote}

\subsection*{You'll want to know: \ta{சீராட்டு}}
\textbf{\ta{சீராட்டு}} — \emph{cīrāṭṭu} — \textquotedblleft{}to pamper; to cherish\textquotedblright{}.

\begin{quote}
\ta{பாட்டி} \ta{பேரனைச்} \ta{சீராட்டி} \ta{வளர்த்தாள்}.
\end{quote}

\begin{quote}
\emph{pāṭṭi pēraṉaic cīrāṭṭi vaḷarttāḷ.}
\end{quote}

\begin{quote}
Grandmother raised her grandson with loving care.
\end{quote}

Past (I …): \ta{சீராட்டினேன்}, \emph{cīrāṭṭiṉēṉ}.

Often paired with \ta{தாலாட்டு} in songs.

\begin{cousinweb}
The origin of \textbf{\ta{சீராட்டு}} (\emph{cīrāṭṭu}) is not settled; it ends in \ta{ஆட்டு}, but our sources do not explain the first part.
\end{cousinweb}

\begin{grammarlens}[title={What you've built}]
One more word for describing what you lived and felt.
\end{grammarlens}

\subsection*{Your turn}
\begin{itemize}
\raggedright
  \item \emph{Say it:} \emph{cīrāṭṭu}
  \item \emph{Say it:} \emph{cīrāṭṭu}, once more
  \item \emph{Say it:} \emph{pāṭṭi pēraṉaic cīrāṭṭi vaḷarttāḷ.} — Grandmother raised her grandson with loving care.
  \item \emph{Recall:} say the Tamil for to adopt, then the Tamil for to coo over, then say \emph{cīrāṭṭu} again
\end{itemize}

\begin{tcolorbox}[breakable,colback=teal!4,colframe=teal!35!black,title={Before you move on}]
What does \ta{சீராட்டு} mean? (\textquotedblleft{}To pamper; to cherish\textquotedblright{}.) Say it once more.
\end{tcolorbox}
\section[\texorpdfstring{\ta{கீழ்ப்படி} (kīḻppaṭi)}{கீழ்ப்படி (kīḻppaṭi)}]{\ta{கீழ்ப்படி} (kīḻppaṭi) — to obey}

\label{lesson:TA-C497-kilppati}



\begin{quote}
Before the new one: say the Tamil for to coo over, then the Tamil for to pamper.
\end{quote}

\subsection*{You'll want to know: \ta{கீழ்ப்படி}}
\textbf{\ta{கீழ்ப்படி}} — \emph{kīḻppaṭi} — \textquotedblleft{}to obey\textquotedblright{}.

\begin{quote}
\ta{குழந்தைகள்} \ta{பெற்றோருக்குக்} \ta{கீழ்ப்படிந்தார்கள்}.
\end{quote}

\begin{quote}
\emph{kuḻantaikaḷ peṟṟōrukkuk kīḻppaṭintārkaḷ.}
\end{quote}

\begin{quote}
The children obeyed their parents.
\end{quote}

Past (I …): \ta{கீழ்ப்படிந்தேன்}, \emph{kīḻppaṭintēṉ}.

\ta{கீழ்} \textquotedblleft{}below\textquotedblright{} + \ta{படி} \textquotedblleft{}to settle\textquotedblright{}. The one obeyed takes -\ta{க்கு}.

\begin{cousinweb}
\textbf{\ta{கீழ்ப்படி}} (\emph{kīḻppaṭi}) is a native compound of \ta{கீழ்} \emph{kīḻ}, \textquotedblleft{}place below\textquotedblright{} (DEDR 1619), and \ta{படி}, \textquotedblleft{}to settle, to submit\textquotedblright{}.
\end{cousinweb}

\begin{grammarlens}[title={What you've built}]
One more word for describing what you lived and felt.
\end{grammarlens}

\subsection*{Your turn}
\begin{itemize}
\raggedright
  \item \emph{Say it:} \emph{kīḻppaṭi}
  \item \emph{Say it:} \emph{kīḻppaṭi}, once more
  \item \emph{Say it:} \emph{kuḻantaikaḷ peṟṟōrukkuk kīḻppaṭintārkaḷ.} — The children obeyed their parents.
  \item \emph{Recall:} say the Tamil for to coo over, then the Tamil for to pamper, then say \emph{kīḻppaṭi} again
\end{itemize}

\begin{tcolorbox}[breakable,colback=teal!4,colframe=teal!35!black,title={Before you move on}]
What does \ta{கீழ்ப்படி} mean? (\textquotedblleft{}To obey\textquotedblright{}.) Say it once more.
\end{tcolorbox}
\section[\texorpdfstring{\ta{பராமரி} (parāmari)}{பராமரி (parāmari)}]{\ta{பராமரி} (parāmari) — to look after; to maintain}

\label{lesson:TA-C497-paramari}



\begin{quote}
Before the new one: say the Tamil for to pamper, then the Tamil for to obey.
\end{quote}

\subsection*{You'll want to know: \ta{பராமரி}}
\textbf{\ta{பராமரி}} — \emph{parāmari} — \textquotedblleft{}to look after; to maintain\textquotedblright{}.

\begin{quote}
\ta{அவள்} \ta{வயதான} \ta{பெற்றோரைப்} \ta{பராமரித்தாள்}.
\end{quote}

\begin{quote}
\emph{avaḷ vayatāṉa peṟṟōraip parāmarittāḷ.}
\end{quote}

\begin{quote}
She looked after her elderly parents.
\end{quote}

Past (I …): \ta{பராமரித்தேன்}, \emph{parāmarittēṉ}.

Of people, houses and machines: \ta{வீட்டைப்} \ta{பராமரி}, to maintain a house.

\begin{cousinweb}
The origin of \textbf{\ta{பராமரி}} (\emph{parāmari}) is not settled. It looks like a Sanskrit loan made into a Tamil verb in -\ta{இ}, perhaps from \emph{parāmarśa}, \textquotedblleft{}seizing; consideration\textquotedblright{} (Monier-Williams), but the link is uncertain.
\end{cousinweb}

\begin{grammarlens}[title={What you've built}]
One more word for describing what you lived and felt.
\end{grammarlens}

\subsection*{Your turn}
\begin{itemize}
\raggedright
  \item \emph{Say it:} \emph{parāmari}
  \item \emph{Say it:} \emph{parāmari}, once more
  \item \emph{Say it:} \emph{avaḷ vayatāṉa peṟṟōraip parāmarittāḷ.} — She looked after her elderly parents.
  \item \emph{Recall:} say the Tamil for to pamper, then the Tamil for to obey, then say \emph{parāmari} again
\end{itemize}

\begin{tcolorbox}[breakable,colback=teal!4,colframe=teal!35!black,title={Before you move on}]
What does \ta{பராமரி} mean? (\textquotedblleft{}To look after; to maintain\textquotedblright{}.) Say it once more.
\end{tcolorbox}
\section[\texorpdfstring{\ta{ஆலோசனை} (ālōcaṉai)}{ஆலோசனை (ālōcaṉai)}]{\ta{ஆலோசனை} (ālōcaṉai) — counsel; consultation}

\label{lesson:TA-C497-alocanai}



\begin{quote}
Before the new one: say the Tamil for to obey, then the Tamil for to look after.
\end{quote}

\subsection*{You'll want to know: \ta{ஆலோசனை}}
\textbf{\ta{ஆலோசனை}} — \emph{ālōcaṉai} — \textquotedblleft{}counsel; consultation\textquotedblright{}.

\begin{quote}
\ta{திருமணம்} \ta{பற்றி} \ta{அப்பா} \ta{தாத்தாவிடம்} \ta{ஆலோசனை} \ta{கேட்டார்}.
\end{quote}

\begin{quote}
\emph{tirumaṇam paṟṟi appā tāttāviṭam ālōcaṉai kēṭṭār.}
\end{quote}

\begin{quote}
Father asked grandfather’s counsel about the wedding.
\end{quote}

\ta{ஆலோசனை} \ta{கேள்}, to ask advice; \ta{ஆலோசனை} \ta{செய்}, to consult.

\begin{cousinweb}
\textbf{\ta{ஆலோசனை}} (\emph{ālōcaṉai}) is borrowed from Sanskrit \emph{ālōcana}, \textquotedblleft{}seeing, perceiving, considering\textquotedblright{} (Monier-Williams).
\end{cousinweb}

\begin{grammarlens}[title={What you've built}]
One more word for describing what you lived and felt.
\end{grammarlens}

\subsection*{Your turn}
\begin{itemize}
\raggedright
  \item \emph{Say it:} \emph{ālōcaṉai}
  \item \emph{Say it:} \emph{ālōcaṉai}, once more
  \item \emph{Say it:} \emph{tirumaṇam paṟṟi appā tāttāviṭam ālōcaṉai kēṭṭār.} — Father asked grandfather’s counsel about the wedding.
  \item \emph{Recall:} say the Tamil for to obey, then the Tamil for to look after, then say \emph{ālōcaṉai} again
\end{itemize}

\begin{tcolorbox}[breakable,colback=teal!4,colframe=teal!35!black,title={Before you move on}]
What does \ta{ஆலோசனை} mean? (\textquotedblleft{}Counsel; consultation\textquotedblright{}.) Say it once more.
\end{tcolorbox}
